% ============================================================
%  CHAPTER I
%  Painting and Relief of Ancient Egypt. Artistic Technical Drawing
%  Source pp. 15–41
% ============================================================
\chapter{Painting and Relief of Ancient Egypt.\newline Artistic Technical Drawing}
\markboth{Chapter I. Painting and Relief of Ancient Egypt}{Chapter I}

\srcpage{15}

As is well known, the mathematically justified foundation for rational methods of
conveying distant regions of perceptual space on a picture plane is the method of
\textit{central projection}, which leads to the system of linear perspective. By analogy,
one may raise the question of the most adequate~--- and equally mathematically
rigorous~--- method of depicting objective space on a plane. The work of mathematicians
of the eighteenth and early nineteenth centuries, in particular that of the French
geometer Gaspard Monge, showed that this method is the \textit{method of orthogonal
projections}. The essence of this method is that the objects of objective space are
projected onto the picture plane by lines perpendicular to the picture plane. If,
moreover, the object is positioned so that its projection onto the picture plane conveys
the most characteristic features of the object, then the resulting projection will allow
these characteristic features to be seen without any geometric distortion. Since real
objects are three-dimensional, a sufficiently complete representation of them can be
obtained by having simultaneously three mutually perpendicular projections~--- a front
view, a side view, and a top view (plan). If a work of art, whose aim is to convey the
geometric appearance of objective space, must consequently be based on the method of
orthogonal projections, then the question naturally arises of how well this method is
able to serve as a foundation for the conveyance of an artistic image.

The method of orthogonal projections in its modern form, adapted for engineering
activity, is of course ill-suited for artistic purposes. First of all, by its rules each object
is depicted three times, in the form of three projections, in different locations on the
picture plane. This makes direct, immediately sensuous perception extremely difficult,
without which no work of art can exist. Consequently, this method could be used in art
on condition that each object is depicted only once.\footnote{%
  Multiple depictions of one figure may be permitted if it is agreed to convey in this
  way a process unfolding in time. Since below we shall be speaking only of spatial
  constructions, this aspect will not be taken into account.}
But in that case another difficulty arises~--- the low informativeness of such a partial
image. Thus the most reasonable approach to solving the arising problem is indicated:
the use of the method of orthogonal projections to obtain a single, most characteristic
depiction of each object, and the increase of the informativeness of this image by various
conditional techniques.\footnote{%
  Such an approach is also characteristic of engineering drawing: if it is possible to give
  a representation of a part by depicting only one of its projections~--- supplemented with
  one or another convention~--- this is always recommended.}

The strict mathematical justification of the method of orthogonal projections by no
means implies that these mathematical regularities must be known to the artist, especially
since the method in question is, in its practical aspect, sufficiently simple and
self-evident.\footnote{%
  One should not underestimate the mathematical knowledge of the Egyptians; they were
  born geometers.}

It is interesting to trace how precisely the ancient Egyptian artists followed the method
of orthogonal projections, and how close the spatial constructions they employed were
to the mathematically optimal.

Let us examine the geometric peculiarities of ancient Egyptian painting and relief
connected with the aspiration to convey to the viewer the geometry of objective space
on a plane, in the following order:

\begin{enumerate}
  \item use of the method of orthogonal projections;
  \item conditional drafting techniques: (a) conditional rotations of planes of
    depiction, (b) cross-sections, (c) scale variation, (d) shifts;
  \item the sign character of images.
\end{enumerate}

\medskip

\noindent\textbf{1. The use of the method of orthogonal projections}, as actually
observed in ancient Egyptian art, leads to a whole series of peculiarities that largely
determine the geometry of ancient Egyptian depiction. First of all, the method in
question recommends a quite definite orientation of the depicted object relative to the
picture plane~--- one at which its most characteristic geometric features are most fully
conveyed. Since in a work of art the object must be shown only once, the method of
orthogonal projections makes it desirable to depict the object in its most characteristic
orientation~--- either from the side, from the front, or from above. This requirement is
not, of course, an absolute law; an object may be depicted from an arbitrary direction,
in foreshortening; but it is a desirable condition, which the ancient Egyptian artist as a
rule follows~--- and it gives ancient Egyptian painting its particular character.

As a rule, for the depiction of human figures or animals the side view is chosen. This
is entirely natural, for the side view is often more informative than the front view or~---
still more~--- the top or rear view. From the front, for example, a standing and a walking
figure would be practically indistinguishable. Analogous considerations apply to the
depiction of animals.\footnote{%
  The geometric basis of this "natural" approach is that when seen from the side, the
  plane of symmetry of the depicted figures is parallel to the picture plane.}
At the same time, slain enemies lying on the ground are shown from above~--- that is,
also from the most characteristic direction.

Another important peculiarity of the depiction of objects using the method of
orthogonal projections is the independence of dimensions on the picture plane from the
distance between the object and that plane. Therefore both a nearby and a distant human
figure will have the same size on the picture plane; if figures of different sizes appear
in the image, they have other, non-projectional causes, which we shall examine later.

\srcpage{16}

When depicting this or that figure, the artist enjoys a certain freedom; but the
depiction of the ground~--- the surface on which these figures stand~--- is subject to
requirements that have a compulsory character. The point is that any objects, people, or
animals can occupy, relative to the top view and two mutually perpendicular horizontal
directions of projection, any position whatsoever~--- but this absolutely cannot be said
of the surface of the ground. The ground surface is visible only in the top view, in plan.
With horizontal directions of projection (front view and side view) the image of the
ground surface reduces to a line; the horizon and the ground surface in the foreground
merge, and figures of people seem to be "standing on the horizon." Therefore the
depiction of the ground surface, first and foremost, testifies to the system of projections
used by the artist. When a drafting projection is used, the person will always stand "on
the horizon"; his feet will never be lower than this line. Even the most modest
\textit{pozem} (ground line) in medieval icons is fundamentally different from the
ancient Egyptian depiction of the ground surface, since the feet of the figures depicted
in the icon are projected onto the image of "earth" and not of "sky."

All this means that in ancient Egyptian painting the ground surface is depicted in the
form of a clear, usually straight horizontal line, which we shall call below the
\textit{baseline}.\footnote{%
  H.~Schäfer in his books (see, for example: H.~Schäfer, \textit{Von ägyptischer Kunst:
  Eine Grundlage}, Wiesbaden, 1963) calls this line \textit{Standlinie}: "the standing
  line." In Russian, such a term gives the introduced concept an overtone of
  immobility, which is undesirable since people frequently run, horses gallop, and so
  on, along this line.}
Almost every time the ancient Egyptian artist depicts people or animals, a baseline is
shown beneath their feet so that they should not appear to "hang in the air."

It should be said that this line acquires meaning only in the system of orthogonal
projections, as the side projection of the surface of the earth. The Egyptians considered
it "firm and impenetrable," like the earth itself. This may be verified by an expressive
illustration from the Book of the Dead, which shows how the soul of the deceased
drinks water from a river in the other world (Fig.\,1). So that the soul might drink the
water, the artist interrupted the depiction of the baseline.

\figblock{1}{Soul of the deceased drinking water in the afterlife. Illustration from the Book of the Dead. 1085--950 BCE.}

The method of depicting the ground surface described here leads to an original solution
of the problem of conveying spatial depth. If it is necessary to convey the comparatively
slight depth of some "layer" of space, one can use the baseline and the single pictorial
sign of depth permissible in technical drawings~--- occlusion (a near object conceals a
more distant one) (see Chapter II). If, however, it is necessary to show deep space, then
the only method of conveying depth is recourse to the plan, the top view; for depicting
near and far objects at the same size~--- without diminution as distance increases~--- makes
a deep "layer" of space indistinguishable from a shallow one. Recourse to the plan also
becomes inevitable when it is necessary to show formations on the ground surface~---
for instance, a pond, a river, and so on~--- all that which with any side projection would
merge with the baseline.

Thus, the method of orthogonal projections, in contrast, for example, to the system of
linear perspective, does not permit the depiction of horizontal surfaces other than in plan.

\srcpage{17}

\figblock{2}{Osiris at a pond with trees and a grapevine. Illustration from the Book of the Dead. 15th century BCE.}

One should particularly emphasize that the flat character of images is an evident
property of every technical drawing in which dimensions are given in their true
proportions, without the apparent diminution of dimensions as distance to the depicted
objects increases.

The foregoing considerations show that art based on the method of orthogonal
projections will possess features that will seem "strange" to a modern person accustomed
to pictures~--- that is, to images of perceptual space.

This strangeness is further heightened for the modern viewer by the circumstance that,
wishing to increase the informativeness of the image, the ancient Egyptian artist was
compelled by the very logic of the adopted method to employ a whole series of
conditional techniques.

\medskip

\noindent\textbf{2a. Conditional rotations of planes of depiction} have the aim of
compensating for the loss of geometric information connected with the need to show
only one projection of an object. In such cases, the transmission of some part of a
depicted figure or object in a conditionally rotated position is permitted. To explain the
use of this technique by the ancient Egyptian artist, one may refer to what first strikes
the eye: the depiction of the human figure. Its unusual appearance is connected with
the fact that while the main direction of projection is the side view (for the legs, body,
and head), the shoulders and eyes are conveyed in the front view~--- that is, in a
conditional rotation. As can be seen from the foregoing, all parts of the body are
depicted so as to show their most characteristic objective features. This is, as it were, a
desire to "spread" the objective three-dimensional geometry onto the two-dimensional
picture plane.

Although this method of depicting the human figure is the most characteristic for
ancient Egyptian art, the artist was not bound by it when other techniques for
conveying objective geometry proved preferable and the conditional rotation turned out
to be unnecessary. Thus, in depicting the human figure the shoulders by no means
always appear in the front view. If a person was depicted whose activity made the
position of his shoulders from the side more expressive, this was always taken into
account. Such departures from the canon usually arose when human labour was being
depicted~--- a ploughman holding a plough with hands drawn close together; a harpist
plucking strings; a sailor climbing a rope; and so on~--- all these personages could lack
the rotation of the shoulder-line. When animals are depicted in the usual side view,
horns are shown from the front in the case of cows, but not antelopes, for which the
side view of the horns is more characteristic.

\srcpage{18}

\figblock{2}{[see caption above; Fig.\,2 caption repeated for in-text reference: Osiris at a pond, Book of the Dead, 15th century BCE]}

If it was necessary to depict a locality with a pond~--- that is, deep space~--- then, as
already noted, the artist had recourse to the possibilities offered by the plan. However,
the trees surrounding the pond, when viewed from above, would become extremely
inexpressive. Therefore they were depicted in a conditional rotation through a right
angle, as a rule with their crowns directed outward relative to the boundaries of the
pond (Fig.\,2), conveying in this way their actual perpendicularity to the shoreline.\footnote{%
  It is sometimes asserted that such depictions "can arise only if the artist mentally
  places himself at the centre of the depicted space"~--- in this case, the pond surrounded
  by trees (B.\,A.~Uspensky, "On the Semiotics of the Icon,"~in: \textit{Studies in Sign
  Systems}, V, Tartu, 1971, pp.~197--198). This categorical assertion cannot be
  accepted. The plan is the most natural method of depicting the earth's surface; it is
  already used, according to ethnographers, by peoples standing at the lowest stages of
  development. Moreover, the plan is by no means necessarily connected with the visual
  perception of space; it is based rather on the experience of human locomotion. It is
  well known that Polynesians can sketch on a sheet of paper the geometry of the layout
  of the islands they visit, which it is inconceivable to see simultaneously from any
  position whatsoever. A plan can have no "centre"; it is equitable everywhere.
  As for the trees being depicted with crowns outward, this is done simply because
  their perpendicularity to the shoreline cannot be shown otherwise. If the artist were
  to turn them in the other direction, they would conceal the image of the pond.
  Moreover, such a depiction is not obligatory. The lower (from the viewer's
  standpoint) row of trees is often shifted downward so as to show the trees with
  crowns pointing upward, without concealing the pond. This, of course, does not
  change the substance~--- the representation of the pond in plan.}
When showing an enclosure for the pharaoh's or nobleman's hunt, the ancient Egyptian
artist conveyed the enclosure in plan, and the animals themselves and the enclosure
fence in side view, in conditional rotation.

This method of depicting deep space is entirely natural: it is used with success today
in tourist sketch-maps, when architectural monuments are shown from the side on a
locality plan. At times this planimetric character of depiction is not so obvious to the
modern viewer. In one ancient Egyptian relief, the transport of the colossal statue of a
nomarch is shown (Fig.\,3).\footnote{%
  \textit{[Fig.\,3 — Transport of a colossal statue of a nomarch.
  Relief. (Cited as Fig.\,3 in the source.)}\,]}
The ropes by which sledges bearing the statue are pulled are shown in plan, while the
people pulling them and the statue itself appear in a conditional side-view rotation;
consequently, the rows of haulers placed one above another are to be understood as
displaced in the "near-to-far" direction.

When human or animal figures in conditionally rotated positions were depicted on a
locality plan, their feet were almost always placed on baselines~--- that is, a local side-
projection image of the ground surface was normally given. These baselines were not
necessarily straight horizontal lines; they could have an "undulating" character, for
example when animals were depicted inside a hunting enclosure.

In depicting objects, the combination in a single image of two mutually perpendicular
projections was permitted (for example, when depicting a bed~--- a side and a top view),
which must also be attributed to the technique of conditional rotations of planes of
projection.

In general, this technique is most widely used in Ancient Egypt.

\srcpage{19}

\figblock{3}{Transport of a colossal statue of a nomarch. Relief (after Schäfer).}

\medskip

\noindent\textbf{2b. Cross-sections} have the aim of increasing the informativeness of
the image. A basket filled with fruits may be shown by the ancient Egyptian painter in
cross-section, so that it may be clear with what exactly it is filled. Showing bird-
catchers who carry their catch in cages, the artist depicts the cages themselves in cross-
section, so that there can be no doubt whatsoever about the contents of the cages. Even
the depiction of a three-story house in cross-section is known, showing the staircases,
the structure of the floors, and many other structural details (Fig.\,4).

\figblock{4}{Cross-section of a three-storey house. New Kingdom (after Schäfer).}

This technique is sufficiently widespread, and its meaning is entirely obvious. This
allows us to forego presenting additional examples.

\medskip

\noindent\textbf{2c. Scale variation.} The ancient Egyptian painter, conveying objective
space on the picture plane, understood perfectly well what additional possibilities
scale variation opens up to him. In ancient Egyptian painting, scale variation is
explained by the aspiration toward greater informativeness, by compositional
considerations, and by the rules for conveying hierarchy.

In connection with the tasks and details of the story, the artist makes warriors
enormously large in comparison with the fortress near which a battle is taking place;
birds sitting on the branches of a tree are often so huge that it is unclear how the
branches can support them (but every feather can be seen and the birds can easily be
identified by species).

\srcpage{20}

\figblock{5}{Procession of servants carrying provisions for the deceased. Relief from a mastaba at Saqqara. c.~2340 BCE.}

The second impulse driving the ancient Egyptian artist toward scale variation was the
desire to improve the composition. Let us cite one example here. In the burial chamber
of a nobleman of the Fifth Dynasty, a procession of servants bringing provisions to the
deceased is shown (Fig.\,5). So that each servant in this procession might be clearly
visible, the large animals which the servants drive~--- and which could conceal them~---
are depicted reduced in size. Thus one of the servants leads an ox that barely reaches
his knee in the depiction; a goose, also carried by this servant, is shown at the same
scale, so that ox and goose turn out to be the same size. In the case examined, the sole
aim of introducing scale variation was the desire to give the entire scene a composition
in which nothing would hinder the beholding of the stately, dignified, ceremonious
procession. Analogous causes impel the ancient Egyptian to neglect strict scale even
when conveying large areas of ground surface in plan.

Quite frequently scale variation has a hierarchical significance: the figure of the
pharaoh is much larger than the figures of the other persons in the same depiction.
Sometimes several gradations may be encountered: largest of all is the figure of the
pharaoh, then (in decreasing order) the figures of noblemen, and smallest of all the
images of common people~--- soldiers, servants, and so on. As regards interacting
personages, the rule of hierarchical enlargement of figures was not usually applied.

\srcpage{21}

In more complex compositions one can observe the simultaneous action of various
causes~--- of different origin~--- leading to scale variation. In one version of the widely
current subject showing the Judgment of Osiris in the afterworld (Fig.\,6), Osiris is
shown on an enlarged scale relative to the other gods. This is undoubtedly a hierarchical
enlargement of scale, since he is the king of the afterworld. But the soul of the deceased
is shown on the same scale as the other gods, although it stands much lower than they
in the hierarchy. This is caused by the fact that the soul which has come to the judgment
is an important personage of the story, and a reduction in its scale would be unjustified.
On the other hand, the 42 gods constituting the tribunal are shown on a reduced scale~---
even relative to the soul of the deceased~--- since they are secondary to the narrative,
and a larger scale for them would have caused compositional difficulties in placing so
large a number of figures in a limited space.

\figblock{6}{Judgment of Osiris in the afterworld. Illustration from the Book of the Dead.}

Thus, scale variation in ancient Egyptian art was a highly characteristic technique; it
was employed for different purposes, and the rules for its use were sufficiently flexible.

\medskip

\noindent\textbf{2d. The shift.} The method of orthogonal projections, when the depicted
human figures are positioned in a "standard" way relative to the picture plane~--- for
example, precisely from the side~--- can lead to great difficulties. Let us explain this
with the following example. There is known a depiction of Pharaoh Akhenaten receiving
a foreign embassy, seated next to his wife (Fig.\,7). The wife's existence can be guessed
from the depiction of part of her arm embracing the pharaoh, and the palm of her other
hand interlaced with the pharaoh's hand. The wife's figure is almost entirely concealed
by the depiction of Akhenaten. Although such a depiction must be regarded as formally
correct, it is far less expressive than another~--- used in more ancient times (Fig.\,8)~---
where the artist introduced a conspicuous convention: he "shifted" the wife's figure
relative to the husband's. As a result of such a conditional shift, the modern viewer
might think that the wife is sitting behind the husband; but the ancient Egyptian viewer,
knowing the customs of the land, understood perfectly that the couple is sitting side by
side (we know this too, since in addition to ancient Egyptian painting, ancient Egyptian
sculpture has survived). It is not surprising that this informationally and artistically
justified method of depiction~--- using a conditional shift~--- continued to be employed after
the unfortunate innovation of the Akhenaten period.

\figblock{7}{Akhenaten with his wife. Early 14th century BCE (after Schäfer).}

\figblock{8}{Depictions of married couples. Left~--- Old Kingdom, right~--- New Kingdom (after Schäfer).}

\srcpage{22}

Conditional shifts allow the showing of what with ordinary methods of depiction would
be hidden from the viewer. This technique finds the most widespread application in
ancient Egyptian art. Depicting a cosmetic vessel stored inside a cup-shaped container
with a lid, the artist found it necessary to show the vessel and the lid separately~---
"floating" above the container~--- and above the vessel, the lid (Fig.\,9). This is a
typical example of shifts of images: in reality the cosmetic vessel is inside the container,
closed by the lid; but without the shift, such a depiction would give the viewer no
information about the contents of the container.

\figblock{9}{Vessel with lid for cosmetics. Old Kingdom.}

From the same considerations, the ancient Egyptian artist, when depicting receptacles
(baskets, pots, vases, and so on), often shows their contents as though "floating" in the
air. The same technique is used when it is necessary to increase information about
people's activities. For example, when relating the toilette of a noble Egyptian lady, the
master constructs the composition so that a chest is visible behind the servant's back
and above it the articles of toilette~--- an indication of the contents of the chest, or a
showing, on an enlarged scale, of the articles that will adorn the court lady.

The technique is also found when craftsmen at work are depicted: one may encounter,
for example, a cobbler near whose head there "floats" an assortment of the various
instruments used in shoemaking, while the cobbler himself is using only one of them at
that moment.

\srcpage{23}

\figblock{3}{[\textit{(See also above.)} Transport of a colossal statue of a nomarch. The haulers are to be understood as displaced in the "near-to-far" direction, not actually arranged in rows.]}

The peculiarities of ancient Egyptian painting examined here~--- the systematic use of
orthogonal projections supplemented with conditional techniques~--- make its geometric
structure very close to engineering drawing. All four main types of conventional
techniques in ancient Egyptian painting: rotations of planes of depiction, cross-sections,
scale variation, and even shifts~--- are today universally accepted and central also in
engineering drawing.\footnote{%
  In modern engineering drawing, rotations of projection planes are a frequently used
  technique; it is permissible to depict one part of a component from one direction and
  another part from a different direction. Consequently, even the ancient Egyptian
  depiction of the human figure is drawing-consistent.
  \par Scale variation occurs in engineering drawing when depicting individual
  components that form part of a single set of documentation, just as in ancient
  Egyptian painting scale variation is applied to different figures and objects in a single
  composition.
  \par Shifts are used today mainly in "everyday" drawing, in the various instruction
  manuals and descriptions printed for buyers of household appliances, to show them
  "what, how, and where is placed" in an assembled appliance. As for cross-sections,
  their use in drawing is well known.}
But then the strict mathematical justification of drawing can be transferred to ancient
Egyptian painting~--- although of course the painters based their work not on
mathematical knowledge but on empirical rules, just as modern artists who follow the
system of linear perspective generally do not know the mathematical subtleties
associated with it.

\srcpage{24}

As paradoxical as it may seem, precisely this aspiration described above~--- to convey
the objective forms of objects without any distortion~--- led (and will always lead) to
images in which, to a greater or lesser degree, the immediate obviousness of ordinary
visual perception is lost. An object depicted in this way must either be transformed in
our consciousness into the true three-dimensional image, or be perceived as a sign. With
regard to ancient Egyptian art, most probably both hold true simultaneously; in any case,
its sign character is indubitable.

\medskip

\noindent\textbf{3. The sign character of ancient Egyptian painting.} To explain this
peculiarity of the art in question, let us cite some examples, beginning with the
depiction of the human figure~--- a depiction that testifies not only to the artist's
aspiration to convey objective geometry, but also to his understanding of the sign
character of this depiction. In conveying the appearance of a walking person, when both
feet placed at a walking stride are visible, both are shown from the side of the big toe.
It is perfectly obvious that such a depiction is absurd; however, it becomes
understandable if we suppose that the artist was depicting not legs as such, but
"signs" of legs. From the sign standpoint, both legs are identical, perform identical
functions, and therefore an identical depiction of them is permissible. To speak of
incapacity on the part of the artist who painted walls or carved a shallow relief is
absurd, the more so as compositions are known in which, for semantic reasons,
"correct" and "unnatural" depictions of legs stand side by side.

\srcpage{25}

In the ancient Egyptian relief cited above (Fig.\,3), the statue being transported is met
by ranks of people walking in step. The involuntary wish to see in this a rudiment of
perspectival depiction is contradicted by the entire geometric structure of ancient
Egyptian art. One can also speak of the depiction here~--- using the method of
orthogonal projections~--- of a kind of rank arranged obliquely relative to the picture
plane; but it is probably more correct to understand such a rank as a sign for "a group
of walking people," obtained by the multiple repetition (with a slight shift) of the figure
of a walking person. The modern viewer perceives the depiction as "marching in step";
but this perception is based on a transference into Ancient Egypt of our present-day
notions. Against this interpretation is the circumstance that when depicting, say, a herd
of cows, the ancient Egyptian artist uses the same technique, and as a result the animals
in his depiction walk "in ranks" and "in step." Consequently, before us is not a proof of
the ancient Egyptian artist's intuitive striving toward realistic picture-making, but on the
contrary rather a testimony to his clear understanding of the sign character of his work.

When depicting a pond, the master uses a series of conventionally geometric "waves"
(Fig.\,2). This is the sign for water; in the same way the water of a river is conveyed;
ships sail on equally depicted water; in the same kind of zigzag pair of lines the water
from a vessel flows conventionally; and so on. Fish and underwater creatures in a
waterway are often depicted on this conventional surface of the water also only as signs
of the inhabitants of the underwater world~--- not with the aim of conveying their actual
life in the water.

The sign character is an inalienable property of ancient Egyptian painting. When a table
with victuals or a tray with jewels is conveyed, all these objects are reproduced in a
"standard" form, one above the other, barely touching each other~--- as they can in fact
never lie. The hands of a person holding a heavily laden tray and of one holding a scroll
of papyrus are conveyed identically, in a position excluding the holding of either the
one or the other.

The examples cited here (and they could easily be multiplied) are sufficient to
understand this peculiarity of ancient Egyptian visual art. One should recall that modern
engineering drawing, as is well known, also widely employs various conventional-sign
images, including some that lead to "impossible" constructions. Just as the ancient
Egyptian artist showed two left legs on the figure of a person walking rightward, the
modern draughtsman in an equally conventional way depicts nuts on two projections
of an assembly drawing of a part, although according to the rules of orthogonal
projections different projections should correspond to different images of the nuts.
He is clearly showing a "sign of a nut" rather than a nut, for otherwise one would have
to accept the existence of "strange" octagonal nuts, just as absurd as a person with two
left legs.

\srcpage{26}

\figblock{7}{[\textit{See above.}] Akhenaten with his wife.}
\figblock{8}{[\textit{See above.}] Depictions of married couples.}

In studying the methods of spatial constructions of the ancient Egyptian artist, one may
use the experience of engineering drawing as a modern analogue that permits a deeper
understanding of the painter~--- and this is precisely what has been done above. Such
an affinity involuntarily suggests the idea of a deep kinship between these two methods
of depicting objective space on a plane.

This affinity is manifest in the approach to the problem, in the methods of depiction
developed, and as a consequence in the geometric properties of ancient Egyptian painting
and modern engineering drawing. Nevertheless, the examination of the raised question
cannot be considered complete without discussing the existing views on the causes of the
so unusual appearance of ancient Egyptian images.

The simplest explanation of the peculiarity of ancient Egyptian painting is to emphasize
its naïve character. Thus, in the very valuable monograph of Schäfer,\footnote{%
  H.~Schäfer, op.~cit.}
which contains a great deal of factual material and many interesting commentaries on
the methods of depicting various objects in ancient Egyptian painting, the author
constantly refers to the closeness of ancient Egyptian images to a child's drawing
(chiefly a pre-school child's). He does not, of course, wish by this comparison to
belittle ancient Egyptian art; he sees the affinity in the approach, not in the degree of
perfection. Nevertheless, it is difficult to conceive of a "naivety" lasting more than
2,000 years alongside the evident high level of development of ancient Egyptian painting.
Moreover, it is worth noting that today no one would explain the methods of engineering
drawing by their closeness to a "naïve" child's drawing (although here, too, there is
much in common). The peculiarities of the conveyance of space in ancient Egyptian
painting must have rational explanations. As for the closeness of many methods of
ancient Egyptian painting to a child's drawing, an entirely different formulation of the
problem is legitimate here. The question to be studied is: what are the causes by which
children~--- especially at younger ages~--- prefer to depict objective rather than perceptual
space? As soon as psychologists and pedagogues provide the corresponding explanations,
the closeness of children's drawing and ancient Egyptian painting will become self-
evident; for ancient Egyptian painting, too, is concerned with objective space
alone.\footnote{%
  Instructive comparisons of ancient Egyptian art~--- whose analysis is of independent
  interest~--- and children's drawing can be found in: I.\,Ya.~Glinskaya, "Pre-perspective
  Methods of Conveying Spatial Information and Some Questions of Teaching Drawing
  to Younger Schoolchildren,"~in: \textit{Artistic Education in Schools}, Leningrad,
  1973, pp.~103--132.}

\srcpage{27}

There exists an explanation of the peculiarities of ancient Egyptian art by its
traditionalism, which allows the "naïve" stage to be pushed back into the depths of time.
It is considered that in the process of the transition from pre-class to class, slave-owning
society~--- when the foundations of religious ideas were taking shape~--- the peculiarities
of the most ancient "naïve" monuments of Egyptian art, which had had a religious
character in the period of the birth of this art, were canonically fixed as obligatory.
The most ancient pictorial canons impeded the development of the artists' creativity and
preserved in art such features as the conveyance of the human figure from the side, with
an unnatural rotation of the shoulder-line; the depiction of a locality as from a bird's-
eye view, with simultaneous showing of trees, people, and animals from the side;
the depiction of objects not actually visible to the artist; and so on. In the process of
lengthy evolution the works of art became more realistic; but the ancient religious
canonical requirements made the final victory of realism impossible.\footnote{%
  See, for example: M.\,E.~Mathieu, \textit{The Art of Ancient Egypt}, Leningrad and
  Moscow, 1961, p.~6 and following.
  This theory sometimes leads to such astonishing utterances about the art of Ancient
  Egypt as: "The artist dared not depict the world as he saw it in reality, that is, vividly
  and directly" (\textit{Drawing: A Textbook for Students of Art and Graphic
  Departments of Pedagogical Institutes}, ed.~A.\,M.~Serov, Moscow, 1975, p.~6).}

\srcpage{28}

Judging by the general character of the reasoning of the authors who hold such a view,
the ancient Egyptian painter would have liked to depict the visible world in the picture,
but the force of tradition held him back from "correct" images. However, before asserting
such a thing, one must understand whether the Egyptian painter wanted to do what is
attributed to him.

Sometimes it is quite correctly noted that ancient Egyptian art was interested not in
the visible appearance of things but in "conceptions" of them; however, no geometric
conclusions are drawn from this thesis.

The imprecision of the usual arguments is obviously connected with the fact that they
fail to take into account the difference between a technical drawing and a picture; in
discussing the drafting methods of the ancient Egyptian artists, the authors of such
logical constructions somehow consider that the artist's aim was a picture and not a
technical drawing.

If the view that reduces the problem to traditionalism is correct, then with the
penetration of "realism" into ancient Egyptian art the canonical methods of depiction
should gradually be overcome and replaced by new ones, and the technical drawing
should gradually approach the picture. If, on the other hand, one considers that ancient
Egyptian art throughout its entire history was "artistic technical drawing" and never
strove to become a picture, then its development must be accompanied not by the dying
out of drafting methods of depiction, but on the contrary by their perfecting.

As is known, in prehistoric times the visual arts and writing were so closely fused that
they are considered as a kind of unity. This pictographic stage was characterized by the
conveyance of messages (for which writing now exists) through the use of a series of
conventional images of various scenes, as well as of individual figures of people, animals,
and objects. With time, these two aspects of pictographic images began to acquire an
independent existence, generating on the one hand hieroglyphic writing and on the other
hand painting. In writing, its sign character came to the fore (although very many
hieroglyphs have the form of images of objects and animals), while in painting its artistic
essence came to the fore (although many sign peculiarities were preserved here). In the
process of the further development of ancient Egyptian culture, this process continued;
as a result, writing increasingly lost the character of images of real objects, acquiring a
purely sign character (the transition first to hieratic and then to demotic
script),\footnote{%
  The transition to alphabetic script is not being emphasized here; this is merely
  a confirmation of the ultimate sign character. Equally, what has been said is not
  contradicted by the simultaneous existence of hieratic script and hieroglyphic
  inscriptions; for today alongside our writing there exists even pictography~--- for
  example, in road signs for drivers.}
while in painting the sign peculiarities were increasingly effaced.

\srcpage{29}

To illustrate the latter, let us cite some examples. Comparing the depictions of a
husband and wife sitting on a bench that belong to the Old and New Kingdoms, Schäfer
noted\footnote{%
  H.~Schäfer, op.~cit., S.~180--181.}
that in the period of the Old Kingdom the artists always~--- in violation of the natural
geometry of overlapping (the wife's knees partially overlap the husband)~--- arranged
things so that the husband would overlap the wife in all parts of the image (Fig.\,8,
left), thereby achieving a sign emphasis on his significance. Indeed, in this case the
wife's legs were concealed not only by the husband but by the bench. The wife then
assumed an impossible pose: sitting to the husband's left, she placed her left hand on
his shoulder. In the period of the New Kingdom the artist decisively preferred
naturalness (Fig.\,8, right), considering that for the sign emphasis of the husband's
priority it fully sufficed to show him in front of the wife.\footnote{%
  The abandonment of the "absurd" position of the wife's left hand on the husband's
  shoulder~--- when she is seated to his left~--- and the transition to the "natural"
  depiction need not necessarily be interpreted as a sign of passage from "worse" to
  "better." In ancient Egyptian painting one can find many such "absurdities." It is
  entirely admissible, however, to regard the majority of them as a consciously applied,
  often productive artistic technique of using geometrically contradictory images~---
  to which Chapter~VIII of the present book is devoted.}

The same course of development is visible in the evolution of the depiction of hunting
scenes in an enclosure. In the period of the Old Kingdom, the pharaoh shoots from a
bow "in general," while in the enclosure animals of different species stand or walk
sedately. This depiction has a clearly sign character; in essence it is merely a
communication to the effect that "the pharaoh is hunting in an enclosure for animals."
In a later time, the hunting pharaoh or nobleman is shown in a dynamic pose; he shoots
not "in general" but at animals; the latter run, fall, struck by arrows, maul one another.
Here the aspiration to plastic informativeness and narrative clearly increases. The images
of animals in superimposed tiers~--- bearing a compositional-sign character~--- are
replaced by a free showing of them on the undulating surface of the ground (observing
the rules of drafting: the enclosure in plan, the fence and animals in a conditional side-
view rotation), and consequently sign-quality gives way to pictorial quality.

\srcpage{30}

The discussed process encompasses even so traditional a domain as the depiction of the
human figure: standard figures are often replaced by more freely rendered ones; figures
with both a right and a left leg appear; figures presented in a three-quarter rotation and
en face images arise. All this testifies to an increase of plastic informativeness and a
softening of the sign character of ancient Egyptian painting and relief.

What is remarkable, however, is that this long-acting and effective tendency had not
the slightest effect on the drafting character of the images. In this sense, the art of the
Old and New Kingdoms is identical.\footnote{%
  This says nothing about the stagnation of ancient Egyptian art. After all, the
  developing European painting from the Renaissance to the impressionists~--- for all
  its diversity~--- can, from a formal-geometric standpoint, be united as painting based
  on the system of linear perspective.}
One may even observe that with time the drafting character of images becomes more
strict, more "scientific." Thus, in hunting scenes in enclosures, when animals were
depicted in superimposed "tiers," each tier contained a depiction of earth and sky.
Later, animals are shown on the plan of the enclosure without multiple showings of sky,
sometimes even without baselines, and the depiction becomes drawing-perfect.

The observed development of the drafting approach is of comparatively small
significance, since the main features of these methods are already well represented in
the art of the Old Kingdom.

The main indicator of the artist's use of drafting methods is the following of the method
of orthogonal projections. This means that all figures and objects~--- or rather, their
elements~--- are shown as they would be seen if one stood directly before them, with the
picture plane perpendicular to the line of sight: that is, if one were, say, to paint the
head while placing one's eye before it at head level, and to paint the soles of the feet
while lowering one's eye to floor level. In this connection, the artist cannot have the
"viewpoint" so important for modern art;\footnote{%
  The concepts of "viewpoint" or "multiple viewpoints" are, as a rule, meaningless if
  the geometry of objective space is depicted. Therefore the use of such concepts in
  analysing ancient Egyptian painting or other kindred methods of depicting space
  should be avoided. The best way to understand this is from the following example:
  a person blind from birth is quite able to scratch on the surface of a board a plan
  of his apartment~--- that is, to give a depiction of objective space on a plane.}
objects of depiction are recorded identically both at the centre of a large composition
and at its edges. Another peculiarity of the image obtained by the method of orthogonal
projections is, as already stated, the fundamental impossibility of depicting the floor,
the surface of the earth, the surface of a table, and other horizontal surfaces, unless one
has recourse to the plan or to the conditional rotation of planes of depiction. For this
reason the baseline~--- indicating that objects are on the ground~--- plays such an
important role in this case. Orthogonal projections also do not permit perspectival
foreshortening. Indeed, the illusion of spatial depth is profoundly alien to art based on
the conveyance of the geometry of objective space~--- when one wishes to convey objective
information, it is inappropriate to have recourse to illusions.

\srcpage{31}

At the same time, the method of orthogonal projections does not forbid foreshortened
images. Their comparatively small number in ancient Egyptian art is connected with the
fact that they diminish the sign character of images. From a drafting-sign standpoint,
it is better to record figures of people, animals, and objects from "standard" directions
(side, front, top), since this simplifies the decoding of images by the viewer. But when
in the New Kingdom the attention of artists was increasingly attracted by plastic
informativeness, images arose that are almost indistinguishable from a realistic picture.
The well-known composition of female musicians confirms what has been said (Fig.\,10).
The group is rendered through the use of orthogonal projection and, according to the
terminology adopted in the present book, is a technical drawing and not a picture. What
tells us this is the fact that the faces of the musicians were painted by the artist as
though placing his eye at the level of those faces, while the soles of their feet he painted
as though lowering his head to floor level. Moreover, the floor itself is not shown; it is
replaced by the baseline. One may further note that the eyes of all four musicians are
shown in the front view~--- but even if this were not the case, following the method of
orthogonal projections would force one to regard such a depiction as a technical drawing
and not a picture.\footnote{%
  In support of the reasonableness of this terminology, one may note that if a
  three-dimensional model of this group of musicians were constructed~--- using dolls,
  say~--- and a modern engineer were invited to make a drawing of the group "from the
  front," his depiction would be geometrically exactly the same as in the ancient
  Egyptian painting. Only the eyes would probably be shown more "correctly"~--- not
  giving front-view images of eyes on faces seen in profile.}
The spatial depth of the depicted space is revealed here by the use of overlapping
(mutual occlusion)~--- the only accepted indicator of depth permissible by the rules of
technical drawing, and entirely appropriate for conveying shallow spatial layers.

\figblock{10}{Group of female musicians. Painting from a Theban tomb. 15th century BCE.}

\srcpage{32}

Thus, the painting of Ancient Egypt was throughout the entire duration of its existence
based on drafting methods: it was indeed "artistic technical drawing." But in that case,
a whole range of aspects of ancient Egyptian art must be considered with this in mind.\footnote{%
  In works devoted to the art of Ancient Egypt, it is of course not at all necessary to
  call works of ancient Egyptian painting "technical drawings." It suffices to stipulate
  the drafting nature of the methods of depiction employed. In this book, the concepts
  of "technical drawing" and "picture" will henceforth be used throughout, since they
  are concise, allow a very clear distinction between the two polar approaches to
  depicting space on a plane, and both these approaches are the subject of study.
  \par In some works devoted to the art of Ancient Egypt, attempts are made to
  discern in the New Kingdom period rudiments of perspectival depiction techniques.
  The authors of such attempts act from the best motives, seeking to show that ancient
  Egyptian painting was developing in the "right direction." This should probably not
  be done: the examples they cite are extremely unconvincing, and the artistic depiction
  of objective space (which excludes perspectival techniques) is just as "legitimate" as
  the depiction of perceptual space.}

First of all, it should be said that the conventions examined above presuppose a good
knowledge of these conventions on the part of both artist and viewer. In this
circumstance lies one of the causes of the traditionalism of ancient Egyptian painting.
Art relying on conventionally drafting, including sign, images had no possibility of
changing conventions so precise in their nature. One should also not forget that so long
as art conveys not perceptual but objective space, it will in any case inevitably contain
geometric conventions, and replacing familiar conventions with other, new ones would
produce nothing but confusion. Moreover, the comparison of the techniques used in
Ancient Egypt with those used today in engineering activity shows their astounding
close kinship. Such a kinship is not, of course, connected with any direct borrowing, but
has as its cause the circumstance that both modern engineering drawing and ancient
Egyptian painting selected, for their purposes, optimal conventions~--- and for that very
reason ones not subject to change within their respective cultures. In light of what has
been said, it appears that the current frequent reduction of the problem of traditionalism
merely to the artist's submission to the rigid requirements of ancient and irrational
religious canons oversimplifies the solution of the problem. Moreover, the theory of
traditionalism does not answer the question of the causes of the peculiarity of ancient
Egyptian painting and relief, merely pushing the origins of this peculiarity back into
the depths of time.

\figblock{11}{Egyptian ships in the land of Punt. Relief from the temple of Hatshepsut. Second half of the 15th century BCE.}

\srcpage{33}

Based on the geometry of the method of orthogonal projections, possessing a "set" of
permitted and~--- as has been shown~--- entirely rational conventions, the ancient Egyptian
painter, by skilfully combining the latter, was able to reproduce sufficiently complex
scenes. To illustrate this, let us turn to the following example.

The master depicted the loading of an ancient Egyptian ship with precious raw materials
and rare plants and animals in the distant land of Punt (Fig.\,11). The ship itself is
shown from the side of its hull; it towers over the horizontal surface of the water and
therefore appears to be standing on the horizon. The water between the ship and the
shore (also horizontal, and therefore also represented only by a baseline, along which
part of the dockers walk) is shown in top view. This is indicated by its depiction, which
is indistinguishable from the depiction of water in a pond (Fig.\,2). On the conventionally
shown top-view surface of the water, sign images of the inhabitants of the underwater
kingdom are provided.\footnote{%
  These images may also be interpreted as a representation of water in "cross-section"
  or as a kind of "transparency." However, their interpretation as sign designations
  seems more plausible; this is indicated by their arrangement in a single horizontal
  line, at equal distances from one another~--- and, chiefly, by showing them from
  characteristic directions, partly from the side and partly from above. The
  underwater part of the ship is also not shown.}
The gangplanks by which dockers board the ship are seen as from a direction along the
stern-to-bow axis; therefore they are depicted exactly from the side~--- and consequently
the dockers walking on the ground, when stepping onto the gangplanks, must perform
a turn of a right angle; or (more probably) they are not initially walking exactly along
the shore, and in making their way from the depth of the shore space toward the
gangplanks, they are not making so sharp a turn. But then the rear figures of the porters
must be taken not "literally" but with a somewhat sign "tinge"~--- as a pictorial
communication to the effect that they are walking to the gangplanks "overland," without
specifying the inessential geometric characteristic of the direction of this motion.

\srcpage{34}

The entire scene is interesting in that it is strictly a technical drawing. Within it there
coexist depictions in three mutually perpendicular standard directions: from the side of
the vessel's hull, from the side of its stern, and from above. With remarkable tact, the
ancient Egyptian artist managed nevertheless to avoid the documentary-technical
dryness, giving a representation of the event that is almost natural even for directly
sensuous visual perception.

Although in the given example the scene shown is perfectly clear at first glance, in more
complex compositions~--- sometimes consisting of several scenes~--- the content becomes
comprehensible only through interpretation, through a story. To be able to give such an
interpretation, one must know the meaning of the accepted conventions~--- one must know
the language of the pictorial work. In contrast to pictures born of the Renaissance,
which can simply be viewed, ancient Egyptian images must also be read. This
circumstance always fundamentally distinguishes a picture from a technical drawing;
not only with regard to ancient Egyptian painting, but also with regard to engineering
drawing, it is customary to speak of the need to "read drawings."

In the well-known scene of a lion hunt in the desert~--- painted on the box of
Tutankhamun~--- the pharaoh and his chariot are greatly enlarged in scale compared with
the accompanying persons~--- warriors, fan-bearers, and other members of the suite.
The lions are rendered at the scale of the pharaoh, to emphasize the difficulty of the
hunt and the bravery of the hunter. At first the impression is created that the lions
randomly fill the entire picture plane before the pharaoh. But for one familiar with the
drafting conventions of ancient Egyptian painting it is clear that before him is a plan
of a section of desert: and, excluding the lions on the main baseline along which the
pharaoh's chariot moves, the two animals nearest to him are lying slain and are
therefore, as required, shown from above, while the other lions~--- still alive~--- appear
in conditional rotation from the side. Above the baseline, under the lions' paws, several
slightly blurred depictions~--- also in side view~--- of the irregularities of the sandy
ground are given, while the pharaoh's suite in three clear tiers everywhere moves along
traditional baselines. These tiers, as usual, are conventional and do not give the actual
distribution of the persons in space; the sign character of the painting asserts itself
again here (Fig.\,12).

\figblock{12}{Scene of a pharaoh's hunt. Painted box from Tutankhamun's tomb. 14th century BCE.}

\srcpage{35}

By an analogous scheme, scenes of military action are constructed. On the other side
of Tutankhamun's box, before the battle chariot of the fighting pharaoh, there is visible
a solid "mêlée" of living and slain enemies filling the entire picture plane; and the
differentiation into the lying (shown in plan) and standing (shown in conditional rotation)
is possible only as a result of careful examination and analysis, the more so as all of them
are rendered on a reduced scale. Unquestionably, the artist did not set himself the aim
of giving a clearly perceptible picture of the battle; his task was to depict the mighty
pharaoh striking "tens of thousands" of enemies who are contemptible before His
Majesty.

Ancient Egyptian painting knows still more "intricate" compositions, but in every one
of them the drafting and sign foundations can be discerned.

In the present book~--- devoted to the comparatively narrow question of spatial
constructions in the visual arts~--- it is inappropriate to analyse in substance the artistic
creativity of the ancient Egyptian painters. However, attention should be drawn to some
aspects of such analysis connected with the problem of spatial constructions.

Ancient Egyptian painting unquestionably deserves the closest attention of art
historians. One cannot but agree with V.\,N.~Lazarev, who wrote as long ago as 1930:
"Since Egyptian drawing on a plane was a phenomenon of art, its attractiveness must
have an aesthetic character."\footnote{%
  V.\,N.~Lazarev, "On the Question of the Egyptians' Depiction of the Human Figure
  on a Plane," \textit{Proceedings of the Section of the History of Art Studies,
  RANIION}, vol.~4, Moscow, 1930, pp.~3--20. In this work, V.\,N.~Lazarev
  particularly emphasizes the high artistic perfection of the ancient Egyptians'
  depiction of the human figure, who managed so harmoniously and naturally to combine
  side and front views of its individual parts.}
In the works devoted to ancient Egyptian art, it is quite rightly said of the rigid canons
by which artists were guided, of how, for example, "standards" of proportional division
of the human figure into parts existed and were strictly regulated both for a standing
and for a seated figure, that the outlines of figures were applied to the walls using a
"grid" for the preliminary marking of the wall, and so on. This may create in some
readers the impression of a certain limitation of ancient Egyptian art. However, the
problem of the canon is not so simple. "The more the art~$\ldots$ of an artist is
constrained by rules, the greater the mastery required~$\ldots$ in order to~$\ldots$
create, without violating these rules and moreover relying on them~$\ldots$ a work that
is original and artistically unrepeatable."\footnote{%
  G.\,M.~Fridlender, \textit{The Poetics of Russian Realism}, Leningrad, 1971, p.~17.}

\srcpage{36}

\figblock{12}{[\textit{See above.}] Scene of the pharaoh's hunt, Tutankhamun's box.}

On an unfinished relief in the tomb of Pharaoh Horemheb in the Valley of the Kings,
a bark of the sun god Ra sailing through the underworld is depicted. The god himself
stands in the naos, and before him and behind him are the figures of the personified
Reason and Magic. These last two figures are merely sketched in paint on the wall
(without any "grid," incidentally). This unfinished relief has preserved traces of the
creative search of the ancient Egyptian artist. On the wall there appears the outline of
two variants (differing in height) of the figure of Reason, and three variants of the
figure of Magic (differing not only in height but also in position relative to the naos).
In this connection, the artist was compelled to consider two variants of the outline of
the stern part of the bark.\footnote{%
  A reproduction of this unfinished relief can be found in: W.~Forman and
  B.~Kischkewitz, \textit{Die altägyptische Zeichnung}, Prague, 1971, Taf.~17.}
Here we have the opportunity to observe directly, as it were, the reflections of an artist
striving for the maximum compositional perfection. The examples cited here and other
analogous ones testify to the fact that ancient Egyptian painters were above all artists
in the highest sense of the word.

The conveyance of the geometry of objective space inevitably leads to a flat image,
since it fundamentally does not permit illusions of spatial depth based on any
perspectival techniques. Conditioned by natural causes, the flat character of images
called for its own methods of enhancing expressiveness, of achieving artistic perfection,
those most appropriate to painting of this type. It is not surprising that, in contrast for
example to the art of the Renaissance, the primary importance was acquired by line,
silhouette, rhythm, symmetry and asymmetry, ornamental quality, decorativeness.

If we speak of the rhythmic structure of ancient Egyptian painting and set aside colour
rhythm (very important in this art), then even purely geometric rhythm alone astounds
by its diversity and expressiveness. Here one may encounter the simplest rhythm, the
rhythm of "ranks" of walking figures~--- for example, in the relief conveying the
transport of the nomarch's statue (Fig.\,3); one can observe a more complex rhythmic
structure, when the rhythm is consciously "disrupted"~--- for example, when showing
a "rank" of cows walking "in step," one of them is depicted with lowered head. Still more
complex will be, for example, the rhythmic structure of warriors running "in single file"
in a relief from the Akhenaten period. Here each warrior differs from the others~---
in clothing, armament, height. At the same time, they are united by the direction of
running, similar poses, equal distances between warriors, and sometimes even by the
fact that the warriors' heads (owing to differences in height) create as it were a smooth
undulating line. These observations could be continued, so effective and varied are the
methods of forming rhythmic structures by the Egyptian artists.

\srcpage{37}

\figblock{13}{Hands of a mistress and servant. Relief on a sarcophagus. c.~2040 BCE.}

The aspiration toward expressive line, toward symmetry combined with asymmetry, may
be illustrated by a detail of a Middle Kingdom relief. Let us note the finish and
refinement of the composition consisting of three hands~--- the resting hand of the
mistress and the hands of the servant doing her hair (Fig.\,13). The wrists of the
servant, shown in conditional rotation, form an astoundingly perfect and symmetrical
composition, which is as it were balanced by the image of the hand raising a cup to the
face. The mistress's hand with the cup contains no elements of symmetry; it is
asymmetrical relative to the pair of the servant's hands; and yet all three hands form
a unity, emphasized by the essentially identical depiction of the three wrists and the
gracefully extended three fingers of each. As with every perfect composition, here it is
inconceivable to change any single element of the depiction without worsening it.

The analysis of more complex compositions~--- including several different scenes (often
displaced in time), individual figures, and ornamental elements~--- also speaks of the
ancient Egyptian artist's excellent understanding of the planimetric principle of his art.
In particular, the planimetric character of the painting and relief permitted the
introduction of hieroglyphic texts as equal-status elements. Of course, this was
facilitated by the fact that hieroglyphs had the same pictorial nature; the planimetric
quality of the Egyptian picture made it possible to artistically fill the "empty spaces"
with hieroglyphs. In doing so, the hieroglyphic texts performed not only an informational
but also a decorative function. Many scholars have already noted that texts executed
with great taste often transformed a painting or relief into a kind of tapestry, the more
so since hieroglyphs were frequently made polychrome.

\srcpage{38}

The necessity of having recourse to drafting conventions and of taking account of the
sign character of painting not so much constrained the artist as it opened new
possibilities to him, of which he made skilful use. The drafting-sign conventions
allowed the creation of compositions inconceivable in modern art.

Take, for example, the battle scene governed by the pharaoh, mentioned above. Modern
art knows many images of military commanders leading their troops into battle. Often
such canvases reduce to a depiction of the commander-in-chief in the foreground, who
with finger or baton "indicates" the small figures of the combatants filling the distant
planes. How much more expressive is the ancient Egyptian composition in which the
pharaoh personally strikes innumerable multitudes of contemptible enemy soldiers. And
how artistically justified (and strictly rationalistically~--- senseless) are the fan-bearers
running alongside the pharaoh's battle chariot, lending the entire scene a symbolic
resonance; while the pair of horses driving the chariot not only trample the fallen
enemies, but gallop "in step" and are caught at a moment when they have adopted a
beautiful, refined pose.

With a freedom inconceivable for painters of the modern era, the ancient Egyptian
painter, for compositional reasons, displaces, rotates, "suspends" in space objects; with
the aim of obtaining an expressive and laconic silhouette, he avoids unnecessary
"realism"~--- as may be seen from the depiction of three hands cited above (Fig.\,13),
which cannot be conveniently performing the depicted actions in the shown position of
the fingers~--- and so on. The art-historical analysis of all these aspects of ancient
Egyptian painting that have barely been sketched here could probably uncover much of
interest. Regrettably, today there do not exist substantial works examining ancient
Egyptian painting as a great phenomenon of art~--- in which not only the questions
sketched above would be examined, but also much else that would allow a just
appreciation of the strength and perfection of ancient Egyptian painting and relief.
These problems still await their investigator.

\srcpage{39}

What impelled the ancient Egyptian artists to choose as the geometric basis of their art
not the picture but the technical drawing? The answer to such a question can be given
outside the bounds by which the content of the present book is constrained, and its
solution will probably require no little time. Nevertheless, some considerations on this
matter may be advanced here. Ancient Egyptian painting and relief had the aim of
communicating objective information to the viewer. Reliefs and paintings communicated
of the deeds of the pharaoh or his dignitaries, of military campaigns, hunts, of various
moments of official activity; ritual scenes and scenes connected with ideas of the afterlife
were provided; and so on. Thus the task of painting was in some measure the impassive
communication of certain objective (or considered objective) information. As for
painting connected with tombs, the figures and objects depicted on the walls were
supposed to serve the deceased in actuality in the other world, and for this reason had
to possess an objective character. Did not such an approach to the visual arts also
stimulate the conveyance of objective geometry? Was this not one of the causes by
which the ancient Egyptian artist preferred to depict objective rather than perceptual
space?

The artistic conveyance of objective space on a plane undoubtedly has its strengths.
Modern painting, having turned to the depiction of perceptual space, was forced to
convey it in a certain sense from a chance viewpoint, with objects seen in various
foreshortened positions. In this case, for each viewpoint and each angle of view, the
depiction of one and the same object turned out to be different. But after all, a table,
for example, does not have only a geometry dependent on the viewpoint and angle of
view~--- and therefore "indefinite"~--- but also an unchanging objective form independent
of them. Precisely such immutable essences the ancient Egyptian considered worthy of
depiction. With this aim, he conveyed the objective geometry of objects; he used colour
not "distorted" by the accidents of lighting but objectively inherent in the objects; and
so on. Modern art would probably have appeared to the ancient Egyptian as an art that
had departed from "realism"~--- a kind of "super-impressionism" attempting to record
the impression of something changeable, and thereby missing what is immutable, most
adequately reflecting the cardinal properties of the world in which man lives and acts.

\srcpage{40}

The geometric foundation of the spatial constructions of ancient Egyptian art consisted
of mathematically grounded drafting methods; these methods were freely transformed
without changing their substance. The ancient Egyptian artists in no way strove toward
the picture (toward the conveyance of perceptual space) but quite consciously preferred
to depict objective space~--- all this speaks of a perfected art, and not of the primitiveness
of pictorial means that prevented ancient Egyptian painting from "growing up" to the
required level.

Free of any geometric shortcomings, ancient Egyptian art developed naturally and
freely. Among arts that took as their foundation the depiction of the geometry of
objective space, ancient Egyptian art is the most integral and complete. In the history
of culture it is fully comparable to the art of the modern period~--- when speaking of
geometric pictorial means. Both these arts show what the consistent aspiration toward
the artistic conveyance on the plane of the image of objective space (Egypt, from the
Old Kingdom to the New) and of perceptual space (European art from the Renaissance
to the impressionists) leads to. These are two extreme cases: in each of them, a
consistent point of view on the methods of conveying space on a plane is maintained
without compromise. Both of these points of view are equally reasonable, mathematically
equally rigorously grounded, and at the same time polar to one another. All other
methods of depicting space (medieval art, the most recent art) will draw from both
sources, as vividly represented by these extreme types of spatial constructions.

\srcpage{41}
